\section{Exhaustive bounded support (unchanged, re-read)}

The enumeration of V3 is unchanged on disk and its digest is unchanged:

\begin{itemize}
  \item states: \textbf{320}; donor-conservativity violations: \textbf{0}; ideal-product
    mismatches: \textbf{0};
  \item minimal one-coordinate closure separations: \textbf{25};
  \item donor-product nonclosure countermodels: \textbf{31};
  \item exact full closure-refinement successes: \textbf{155};
  \item proper-subset closure-refinement failures: \textbf{1,055};
  \item heterogeneous composition successes under exact bridge: \textbf{25};
  \item bridge-mismatch composition countermodels: \textbf{25};
  \item canonical rows SHA-256:\\
    \idt{25f40385714adb15bca298a8cfd2b7fe2b28c96bfe462f6b60583be8f735b95f}.
\end{itemize}

What V4 changes is the reading. The 155 and 1,055 are a five-way donor
multiplication of 31 and 211: the donor loop enters neither the closure vector
nor the repair set. Neither does anything else under it --- 320 states is the
64-point (native verdict, closure vector) space enumerated once per family, the
25 minimal separations are 5 counted five times, and the 25 successes and 25
bridge countermodels are one of each counted once per ordered donor pair. Only
the 31 nonclosure countermodels are 31 distinct facts, and the \idt{donor\_axis}
block of the result artifact carries the same table. The composition pair is
discussed in \S20. The counts remain
correct and remain reproducible; they are support for a bounded instance, and
the general claims are now carried by \S\S18--19 instead.
